\documentclass[11pt]{article}
\usepackage[a4paper,margin=28mm]{geometry}
\usepackage{amsmath,amssymb,amsthm}
\usepackage[colorlinks=true,linkcolor=blue,urlcolor=blue,citecolor=blue]{hyperref}
\newtheorem*{theorem*}{Theorem}
\newtheorem*{corollary*}{Corollary}
\newenvironment{statement}[1]{\par\medskip\noindent\phantomsection\textbf{#1}\ }{\par\medskip}
\setlength{\emergencystretch}{3em}
\newcommand{\Q}{\mathbb Q}
\newcommand{\Z}{\mathbb Z}
\newcommand{\F}{\mathbb F}
\DeclareMathOperator{\End}{End}
\DeclareMathOperator{\Gal}{Gal}
\DeclareMathOperator{\GL}{GL}
\DeclareMathOperator{\PGL}{PGL}
\DeclareMathOperator{\Hom}{Hom}
\DeclareMathOperator{\Res}{Res}
\DeclareMathOperator{\Ind}{Ind}
\DeclareMathOperator{\tr}{tr}
\DeclareMathOperator{\disc}{disc}
\title{Rational \(5\)-torsion and the \(2\)-rank of geometrically split Jacobians}
\author{}
\date{11 November 2026}

\begin{document}
\maketitle

\begin{theorem*}
Let \(A/\Q\) be a principally polarized abelian surface that is simple over \(\Q\) and
geometrically split. If \(A(\Q)\) has a point of order \(5\), then \(A(\Q)[2]\cong(\Z/2)^r\) with
\(r\le2\).
\end{theorem*}

\begin{corollary*}
No genus-\(2\) Jacobian over \(\Q\) that is geometrically split and simple over \(\Q\) has a
rational torsion subgroup containing \((\Z/2)^3\times\Z/5\). In particular
\href{https://genus-2-torsion.github.io/group.html?g=2.2.10}{\([2,2,10]\)},
\href{https://genus-2-torsion.github.io/group.html?g=2.2.20}{\([2,2,20]\)} and
\href{https://genus-2-torsion.github.io/group.html?g=2.2.2.10}{\([2,2,2,10]\)} do not occur as the
rational torsion subgroup of such a Jacobian. Together with the
\href{https://genus-2-torsion.github.io/proofs/qsplit-two-rank-obstruction.html}{\(2\)-rank
obstruction for Jacobians split over \(\Q\)}, no geometrically split genus-\(2\) Jacobian over
\(\Q\) has rational torsion subgroup \([2,2,2,10]\).
\end{corollary*}

Here ``geometrically split'' means that \(A_{\bar\Q}\) is isogenous to a product of two elliptic
curves. We impose no restriction on the degree of this isogeny, on the field of definition of the
geometric endomorphisms or on the integral endomorphism ring. The principal polarization is used
in one place only, \hyperref[lemma-full]{Lemma~4}. Along the way we prove a statement that needs
neither a point of order \(5\) nor a polarization:

\begin{statement}{Proposition A.}\label{prop-A}
Let \(A/\Q\) be an abelian surface such that \(A_{\bar\Q}\) is isogenous to the square of an
elliptic curve without complex multiplication, and such that the field of definition of the
geometric endomorphisms of \(A\) is a cyclic quartic extension of \(\Q\). Then
\(A(\Q)[2^\infty]\) has order at most \(4\).
\end{statement}

\emph{Outline.} We follow the case distinction of the
\href{https://genus-2-torsion.github.io/proofs/qsimple-2214-obstruction.html}{proof for
\([2,2,14]\)}. The CM case follows from a point count at \(3\): no abelian surface over \(\F_3\)
has exactly \(40\) points (Section~\ref{cm}). If \(A_{\bar\Q}\) is the square of a curve without
CM, the residual-representation lemmas of the
\href{https://genus-2-torsion.github.io/proofs/qsimple-63-obstruction.html}{proof for order 63},
applied at the prime \(5\), leave only the endomorphism fields that are cyclic of degree \(2\) or
\(4\) (Section~\ref{gamma}). In the cyclic quartic case the determinant character of order \(4\)
forces \(\#\widetilde A(\F_p)\not\equiv0\pmod 8\) at a positive density of primes, which gives
Proposition~A (Section~\ref{quartic}). In the remaining cases \(A\) is isogenous to the Weil
restriction of an elliptic curve \(E\) over a quadratic field \(K\) with \(E^\sigma\neq E\). A
principal polarization and rational \(2\)-rank \(3\) produce such an \(E\) with full
\(K\)-rational \(2\)-torsion, and the point of order \(5\) then gives a \(K\)-point on the modular
curve \(X_1(2,10)\), an elliptic curve of rank \(0\) over \(\Q\). At primes inert in \(K\),
\(\#\widetilde A(\F_p)=\#\widetilde E(\F_{p^2})\) must be divisible by \(8\). By Chebotarev's
theorem this forces certain norms to be rational squares. A quadratic point with rational
\(x\)-coordinate comes from an elliptic curve over \(\Q\), which contradicts simplicity over
\(\Q\); at every other quadratic point that is not a cusp, these norms are not all squares in the
required way (Section~\ref{quadratic}).

\section{Conventions and inputs}\label{inputs}

All number fields are subfields of a fixed algebraic closure \(\bar\Q\). For a number field \(k\)
we write \(G_k=\Gal(\bar\Q/k)\). For a prime \(\ell\), \(\chi_\ell\) is the \(\ell\)-adic
cyclotomic character, \(T_\ell\) is the Tate module and \(V_\ell=T_\ell\otimes\Q_\ell\). For an
abelian variety \(X\), \(\End^0=\End\otimes\Q\) and \(\Hom^0=\Hom\otimes\Q\). Conjugate varieties
and morphisms are as in \cite[Section~1]{P2214}: for \(g\in G_{\Q}\), \(Y^g\) is the conjugate of \(Y\)
and \(g:Y(\bar\Q)\to Y^g(\bar\Q)\) the induced bijection. For a quadratic field \(K\) we write
\(\sigma\) for its nontrivial automorphism, and \(\tilde\sigma\in G_{\Q}\) for a lift of
\(\sigma\). We write \(\Q^{\times2}\) for the nonzero rational squares.

\begin{itemize}
\item \textbf{(I1) Reduction} \cite{ST}, as stated in \cite[Section~1]{P63} and \cite[(I2)]{P2214}:
let \(K'\) be a finite extension of \(\Q_p\) with residue field \(k\), \(A\) an abelian variety
over \(K'\) and \(\ell\neq p\) a prime. (a) Good reduction is detected by the inertia action on
\(T_\ell(A)\) \cite[\S1, Theorem~1 and Corollary~2]{ST}; in particular it is preserved by
\(K'\)-isogenies. (b) If \(A\) has good reduction \(\widetilde A\) and \(m\) is prime to \(p\),
then \(A(K')[m]\) injects into \(\widetilde A(k)\) \cite[Lemma~2]{ST}, and
\(T_\ell(A)\cong T_\ell(\widetilde A)\) compatibly with Frobenius. (c) If \(A\) has potentially good
reduction, then it has good reduction over a finite extension of \(K'\) with residue field \(k\)
\cite[\S2]{ST}. (d) (Weil.) For an abelian variety \(B\) of dimension \(g\) over \(\F_q\),
\(\#B(\F_q)=P(1)\), where \(P\) is the characteristic polynomial of the Frobenius endomorphism on
\(V_\ell(B)\), and the complex roots of \(P\) have absolute value \(\sqrt q\); hence
\(\#B(\F_q)\le(1+\sqrt q)^{2g}\). By (b) and (d), if an abelian variety \(X\) over a number field
\(k\) has good reduction \(\widetilde X\) at a prime \(\mathfrak p\nmid\ell\), then
\(\#\widetilde X(O_k/\mathfrak p)=\det\bigl(1-\mathrm{Frob}_{\mathfrak p}\mid V_\ell(X)\bigr)\)
for an arithmetic Frobenius element \(\mathrm{Frob}_{\mathfrak p}\in G_k\) at \(\mathfrak p\).

\item \textbf{(I2) The case distinction} \cite[(I3)]{P2214}. A geometrically split abelian surface
\(A/\Q\) is of exactly one of the types \textbf{(N)} (\(A_{\bar\Q}\) is isogenous to \(E_1\times
E_2\) with \(E_1\), \(E_2\) not isogenous), \textbf{(CM)} (\(A_{\bar\Q}\) is isogenous to \(E^2\)
with \(E\) a CM curve) and \textbf{(S)} (\(A_{\bar\Q}\) is isogenous to \(E^2\) with \(E\) without
CM). In case (S), \(D=\End^0(A_{\bar\Q})\cong M_2(\Q)\); \(L\) is the smallest field over which
all geometric endomorphisms are defined; \(\Gamma=\Gal(L/\Q)\) embeds in \(\PGL_2(\Q)\), the
element \(g\) acting on \(D\) by conjugation with \(R_g\in D^\times\); and
\(\End^0_{\Q}(A)=D^\Gamma\). The group \(\Gamma\) is one of \(C_1\), \(C_2\), \(C_3\), \(C_4\),
\(C_6\), \(V_4\), \(D_6\), \(D_8\), \(D_{12}\), and \(\Gamma\neq C_1\) if \(A\) is simple over \(\Q\)
\cite[Section~5]{P63}. An element
\(r\in\Gamma\) of order \(n\in\{3,4,6\}\) has a unique representative \(R\) with \(\tr R=1\), and
\(R^2-R+1/c=0\) with \(c=1,2,3\) respectively.

\item \textbf{(I3)} \cite[Lemma~2]{P63} Let \(X\) be an abelian surface over a number field \(k\),
\(E\subset X\) an elliptic curve defined over \(k\), and \(Q\in X(k)\) of prime order \(q\). Then
either \(Q\in E(k)\), or the image of \(Q\) in \((X/E)(k)\) has order \(q\).

\item \textbf{(I4)} \cite[Lemma~3]{P63} In case (S), every elliptic curve in \(A_{\bar\Q}\) is
defined over \(L\), and any two of them are \(L\)-isogenous.

\item \textbf{(I5) Residual representations} \cite[Section~6]{P63}. Let \(p\) be an odd prime and \(A\)
of type (S). Lemma~5 there (a polarization of \(A\) over \(\Q\), of any degree) is used by
Lemma~6 (\emph{unramified order}) and Lemma~7 (\emph{split blocks}); both conclude
\(\Q(\zeta_p)\subseteq L\) when \(A(\Q)[p]\neq0\) and a \(\Gamma\)-stable \(\Z_p\)-order
\(\mathcal B\subset D\otimes\Q_p\) of the stated kind exists.

\item \textbf{(I6) The determinant} \cite[Lemma~3]{P2214}. Let \(A\) be of type (S) with
\(\Gamma=\langle r\rangle\) cyclic of order \(n\in\{3,4,6\}\), and \(R\) as in (I2). Then
\(F=\Q[R]\) is a field, because \(X^2-X+1/c\) has negative discriminant (\(F\cong\Q(i)\) for
\(n=4\)); as \(R\) is not a scalar, its centralizer in \(D\) is \(F\), so
\(F=D^\Gamma=\End^0_{\Q}(A)\); and \(N(R)=\det R=1/c>0\). The statement and proof of
\cite[Lemma~3]{P2214} are written for \(n\in\{3,6\}\), but they use only these facts and that the
roots of unity in \(F\) have norm \(1\); they apply verbatim to \(n=4\), \(c=2\). For an elliptic
curve \(E\subset A_L\) and \(H=\Hom^0(E_{\bar\Q},A_{\bar\Q})\), a one-dimensional \(F\)-vector space,
evaluation is an \(F\)-linear isomorphism \(V_\ell(E)\otimes_{\Q}H\to V_\ell(A)\). Under it each
\(g\in G_{\Q}\) acts as \(B_g\otimes\gamma_g\) with \(B_g\in\GL(V_\ell(E))\) and
\(\gamma_g\in F^\times\), and
\[\varepsilon(g):=\gamma_g^2/\deg(u_g)\]
for a quasi-isogeny \(u_g:E\to E^g\) of degree \(\deg(u_g)\in\Q_{>0}\) is a character
\(\varepsilon:G_{\Q}\to\mu(F)\) that factors through \(\Gamma\), with \(\varepsilon(r)=R/\bar R\)
a primitive \(n\)-th root of unity. The \(F\otimes\Q_\ell\)-determinant of \(g\) is
\(\chi_\ell(g)\varepsilon(g)\).

\item \textbf{(I7) Weil restriction} \cite[\S1]{Mi}. Let \(K\) be a quadratic field and \(E\) an
elliptic curve over \(K\). The Weil restriction \(\Res_{K/\Q}E\) is an abelian surface over \(\Q\)
with \((\Res_{K/\Q}E)(\Q)=E(K)\) and \((\Res_{K/\Q}E)_K\cong E\times E^\sigma\). Its
\(\bar\Q\)-points are the pairs \((x,y)\) with \(x\in E(\bar\Q)\) and \(y\in E^\sigma(\bar\Q)\);
an element \(g\in G_K\) maps \((x,y)\) to \((gx,gy)\), and an element \(g\notin G_K\) maps it to
\((gy,gx)\). Accordingly \(V_\ell(\Res_{K/\Q}E)\cong\Ind_{G_K}^{G_{\Q}}V_\ell(E)\). If \(E\) is
the base change of an elliptic curve \(E_0\) over \(\Q\), then \(\Res_{K/\Q}E\) is isogenous over
\(\Q\) to \(E_0\times E_0^{(K)}\), where \(E_0^{(K)}\) is the quadratic twist.

\item \textbf{(I8) Halving} \cite[Theorem~4.2]{Kn}. Let \(k\) be a field of characteristic
different from \(2\) and \(3\), and \(E:y^2=(x-e_1)(x-e_2)(x-e_3)\) with distinct \(e_i\in k\). A point
\((x_0,y_0)\in E(k)\) lies in \(2E(k)\) if and only if \(x_0-e_1\), \(x_0-e_2\) and \(x_0-e_3\)
are squares in \(k\).

\item \textbf{(I9) Kubert's normal form} \cite[Table~3, the entry for \(\Z/10\), with
\(f=t\)]{Ku}. For \(t\) in a field \(k\) of characteristic \(0\) put
\[c=-\frac{t(t-1)(2t-1)}{t^2-3t+1},\qquad b=\frac{t^3(t-1)(2t-1)}{(t^2-3t+1)^2},\qquad
E_t:\ y^2+(1-c)xy-by=x^3-bx^2 .\]
If \(E_t\) is an elliptic curve, \((0,0)\) has order \(10\), and every pair \((E',P)\) of an
elliptic curve over \(k\) and a point \(P\in E'(k)\) of order \(10\) is isomorphic over \(k\) to
\((E_t,(0,0))\) for some such \(t\in k\). Kubert works over \(\Q\), but his derivation from the
Tate normal form \(y^2+(1-c)xy-by=x^3-bx^2\) with \(P=(0,0)\) uses only field operations and is
valid over every field of characteristic \(0\); the parameter is recovered as \(t=c^2/(b-c)\).
\end{itemize}

\section{The CM case: a point count at \(3\)}\label{cm}

\begin{statement}{Lemma 1.}\label{lemma-f3}
No abelian surface over \(\F_3\) has exactly \(40\) points.
\end{statement}

\emph{Proof.} Let \(B\) be an abelian surface over \(\F_3\). By Weil, the characteristic
polynomial of Frobenius is \((x^2-\beta_1x+3)(x^2-\beta_2x+3)\) with \(\beta_1,\beta_2\) real,
\(|\beta_i|\le2\sqrt3\), and \(h(y)=(y-\beta_1)(y-\beta_2)=y^2-sy+t\) has integer coefficients.
Then \(\#B(\F_3)=(4-\beta_1)(4-\beta_2)=h(4)=16-4s+t\). If \(h(4)=40\), then \(t=24+4s\). The roots
are real, so \(s^2\ge4t=96+16s\), that is \((s-8)^2\ge160\), so \(s\le-5\) or \(s\ge21\); and
\(|s|=|\beta_1+\beta_2|\le4\sqrt3<7\). Hence \(s\in\{-5,-6\}\) and \(h=(y+1)(y+4)\) or
\(h=y(y+6)\), with the root \(-4\) or \(-6\) outside \([-2\sqrt3,2\sqrt3]\). \qed

\medskip
The certificate (Section~\ref{computations}, part A) lists all values of \(h(4)\); the largest is
\(49\).

\medskip
\textbf{The CM case.} Let \(A\) be of type (CM) with a point \(P_5\in A(\Q)\) of order \(5\), and
suppose \(A(\Q)[2]\cong(\Z/2)^r\). As in \cite[Section~3]{P2214}, \(A\) has potentially good reduction
at \(3\), and by (I1)(c) good reduction \(\widetilde A\) over a finite extension of \(\Q_3\) with
residue field \(\F_3\). The subgroup of \(A(\Q)\) generated by \(P_5\) and \(A(\Q)[2]\) has order
\(5\cdot2^r\), prime to \(3\), so by (I1)(b) \(5\cdot2^r\) divides \(\#\widetilde A(\F_3)\),
which is at most \((1+\sqrt3)^4<56\) by (I1)(d). If \(r\ge3\), this forces \(\#\widetilde A(\F_3)=40\), contradicting
\hyperref[lemma-f3]{Lemma~1}. So \(r\le2\). This case uses neither simplicity over \(\Q\) nor a
polarization.

\section{Type (S): the endomorphism field is cyclic of order \(2\) or \(4\)}\label{gamma}

\begin{statement}{Proposition 2.}\label{prop-gamma}
Let \(A\) be of type (S) and simple over \(\Q\), with \(A(\Q)[5]\neq0\). Then \(\Gamma\) is
cyclic of order \(2\) or \(4\).
\end{statement}

\emph{Proof.} We follow \cite[Section~7]{P63} with the prime \(7\) replaced by \(5\), using (I5) with
\(p=5\). The field \(\Q(\zeta_5)\) is cyclic of degree \(4\) over \(\Q\), so if
\(\Q(\zeta_5)\subseteq L\), then \(\Gamma\) has a quotient isomorphic to \(C_4\). Modulo \(5\) the
discriminants \(-3\), \(-1\), \(-\tfrac13\) of the polynomials \(X^2-X+\tfrac1c\) for rotations
of order \(3\), \(4\), \(6\) are \(2\), \(4=2^2\) and \(3\); the nonzero squares modulo \(5\) are
\(1\) and \(4\). By (I2) and simplicity, \(\Gamma\neq C_1\).
\begin{enumerate}
\item \emph{\(\Gamma\in\{C_3,C_6,D_6,D_{12}\}\).} Let \(r\) be a rotation of order \(3\) or \(6\)
generating the rotation subgroup and \(R\) its trace-one representative, with \(c\in\{1,3\}\) a
\(5\)-adic unit. Every element of \(\Gamma\) maps \(R\) to \(R\) or \(1-R\) \cite[Section~7]{P63}, so
\(\mathcal B=\Z_5[R]\) is a \(\Z_5\)-order preserved by \(\Gamma\). Its discriminant is a
non-square unit modulo \(5\), so \(\mathcal B\) is the ring of integers of the unramified
quadratic extension of \(\Q_5\); \(R\in\mathcal B^\times\) because \(\det R=1/c\) is a unit, and
its reduction is a root of an irreducible quadratic over \(\F_5\), hence not in \(\F_5\). By
\cite[Lemma~6]{P63}, \(\Q(\zeta_5)\subseteq L\). But \(C_3\), \(C_6\) and the abelianizations \(C_2\)
and \(C_2\times C_2\) of \(D_6\) and \(D_{12}\) have no quotient \(C_4\).

\item \emph{\(\Gamma=D_8\).} Let \(r\) be a rotation of order \(4\) and \(R\) its trace-one
representative, \(R^2-R+\tfrac12=0\). Again \(\mathcal B=\Z_5[R]\) is preserved by \(\Gamma\). Its
discriminant \(-1\equiv2^2\) is a nonzero square modulo \(5\), so by Hensel's lemma the roots
\(\lambda_1,\lambda_2\) of \(X^2-X+\tfrac12\) lie in \(\Z_5^\times\) with \(\lambda_1-\lambda_2\)
a unit, and \(\mathcal B\cong\Z_5\times\Z_5\) with idempotents
\(e_1=(R-\lambda_2)/(\lambda_1-\lambda_2)\) and \(e_2=(R-\lambda_1)/(\lambda_2-\lambda_1)\). The
rotations fix \(R\), hence \(e_1\) and \(e_2\); the reflections map \(R\) to \(1-R\) and, since
\(\lambda_1+\lambda_2=1\), exchange \(e_1\) and \(e_2\). The stabilizer of \(e_1\) is the rotation
subgroup, and for \(X\in e_1(D\otimes\Q_5)e_2\) we have \(RXR^{-1}=(\lambda_1/\lambda_2)X\), so
the character \(\eta\) of \cite[Lemma~7]{P63} maps a lift of \(r\) to \(\lambda_1/\lambda_2\), the
ratio of the eigenvalues of \(R\), a primitive fourth root of unity. As \(5\nmid4\), its reduction
\(\bar\eta\) has order \(4\). By \cite[Lemma~7]{P63}, \(\Q(\zeta_5)\subseteq L\); but the
abelianization of \(D_8\) is \(C_2\times C_2\).

\item \emph{\(\Gamma=V_4\).} The argument of \cite[Section~7, case \(V_4\)]{P63} applies verbatim with
\(7\) replaced by \(5\): if \(a=R^2\) (for a suitable \(r\) and trace-zero representative \(R\)
with \(\det R\in\Z_5^\times\)) is not a square modulo \(5\), \cite[Lemma~6]{P63} applies; if it is, a
reflection exchanges the idempotents \((1\pm R/b)/2\), \(\bar\eta\) is the nontrivial quadratic
character, and \cite[Lemma~7]{P63} applies. In both cases \(\Q(\zeta_5)\subseteq L\), which is
impossible because \([L:\Q]=4\) and \(V_4\) is not cyclic.
\end{enumerate}
The remaining groups are \(C_2\) and \(C_4\). \qed

\section{Type (S) with a cyclic quartic endomorphism field}\label{quartic}

\emph{Proof of \hyperref[prop-A]{Proposition A}.} Here \(A\) is of type (S) and
\(\Gamma=\langle r\rangle\cong C_4\). By (I6) with \(n=4\), \(F=\End^0_{\Q}(A)\cong\Q(i)\),
\(G_{\Q}\) acts \(F\otimes\Q_\ell\)-linearly on the free rank-\(2\) module \(V_\ell(A)\), and
\(\varepsilon(r)=R/\bar R\in\{i,-i\}\). Take \(\ell=2\). Then \(F\otimes\Q_2=\Q_2(i)\) is a field
with ring of integers \(\Z_2[i]\) and uniformizer \(\pi=1+i\).

Let \(p\) be an odd prime of good reduction for \(A\), unramified in \(L\), whose arithmetic
Frobenius \(g=\mathrm{Frob}_p\in G_{\Q}\) satisfies \(\varepsilon(g)=\pm i\). By (I6), \(g\) acts as \(B\otimes\gamma\)
with \(\gamma\in F^\times\) and \(\gamma^2=\varepsilon(g)\,\delta\), \(\delta=\deg(u_g)>0\).
Writing \(\gamma=x+yi\) with \(x,y\in\Q\), the equations \(x^2-y^2=0\) and \(2xy=\pm\delta\) give
\(\gamma=x(1\pm i)\). Hence the trace of \(g\) over \(\Q_2(i)\) is
\[a=\tr(B)\,\gamma\in(1\pm i)\,\Q_2 .\]
The roots of the characteristic polynomial \(X^2-aX+b\) of \(g\) over \(\Q_2(i)\) are eigenvalues
of Frobenius, hence algebraic integers, so \(a\in\Z_2[i]\); writing \(a=(1\pm i)u\) with
\(u\in\Q_2\), the valuation \(v_\pi(a)=1+2v_2(u)\ge0\) gives \(u\in\Z_2\). By (I6),
\(b=\chi_2(g)\varepsilon(g)=\pm ip\). With (I1),
\[\#\widetilde A(\F_p)=\det\nolimits_{\Q_2}(1-g)=N_{\Q_2(i)/\Q_2}(1-a+b)
=N\bigl((1-u)\pm i(p-u)\bigr)=(1-u)^2+(p-u)^2 .\]
Suppose \(p\equiv3\pmod4\). If \(u\in2\Z_2\), both \(1-u\) and \(p-u\) are units and the sum is
\(\equiv2\pmod8\). If \(u\in1+2\Z_2\), write \(1-u=2v\) and \(p-u=2w\); then \(w-v=(p-1)/2\) is
odd, so \(v^2+w^2\) is a unit and the sum is \(\equiv4\pmod8\). In both cases
\(8\nmid\#\widetilde A(\F_p)\).

There are infinitely many such primes. If \(\Q(i)\subseteq L\), then \(\Q(i)\) is the quadratic
subfield of \(L\), and every prime whose Frobenius generates \(\Gal(L/\Q)\) is inert in \(\Q(i)\),
that is \(p\equiv3\pmod4\). (This case does not occur: the argument of \cite[Section~9.1]{P63},
which uses only that every element of \(\Gamma\) is the class of a power \(R^k\) with
\(\det R^k>0\), shows that \(L\) is totally real. We do not need this.) Otherwise \(L\cap\Q(i)=\Q\), \(\Gal(L(i)/\Q)\cong\Gal(L/\Q)\times
\Gal(\Q(i)/\Q)\), and Chebotarev's theorem gives infinitely many primes whose Frobenius generates
\(\Gal(L/\Q)\) and which are \(\equiv3\pmod4\). For each such \(p\), \(A(\Q)[2^\infty]\) injects
into \(\widetilde A(\F_p)\) by (I1)(b), so its order divides \(4\). \qed

\medskip
Proposition~A uses no point of order \(5\), no simplicity over \(\Q\) (a surface with
\(\Gamma\cong C_4\) is simple over \(\Q\), because \(\End^0_{\Q}(A)\) is a field) and no
polarization.

\section{The quadratic case}\label{quadratic}

In this section \(A\) is simple over \(\Q\), geometrically split, and either of type (N) or of
type (S) with \(\Gamma\cong C_2\). In type (N) let \(K\) be the quadratic field over which the
two elliptic curves of \(A_{\bar\Q}\) are defined (they are exchanged by \(\sigma\)); in type (S)
put \(K=L\). Every elliptic curve \(E\subset A_{\bar\Q}\) that we consider is defined over \(K\),
and \(E^\sigma\neq E\), since otherwise \(E\) would be stable under \(G_{\Q}\) and defined over
\(\Q\) \cite[Section~2]{P2214}, contradicting the simplicity of \(A\) over \(\Q\).

\begin{statement}{Lemma 4 (a curve with full \(2\)-torsion).}\label{lemma-full}
Suppose that \(A\) is principally polarized and that \(A(\Q)[2]\) has rank at least \(3\). Then
there is an elliptic curve \(E_1\subset A_K\), defined over \(K\), with \(E_1^\sigma\neq E_1\) and
\(E_1[2]\subseteq E_1(K)\).
\end{statement}

\emph{Proof.} Let \(\lambda\) be the principal polarization and \(e\) the Weil pairing on
\(A[2]\) attached to it, \(e(z,z')=e_2^A(z,\lambda z')\), where \(e_2^A:A[2]\times A^\vee[2]\to\mu_2\)
is the canonical pairing; \(e\) is a nondegenerate alternating \(G_{\Q}\)-invariant pairing. If \(A(\Q)[2]=A[2]\), any elliptic curve \(E\subset A_K\) defined over \(K\) works. So
let \(W=A(\Q)[2]\) have dimension \(3\). Each \(g\in G_{\Q}\) acts trivially on \(W\) and on the
line \(A[2]/W\), so \(g=1+N_g\) with \(N_g(A[2])\subseteq W\) and \(N_g(W)=0\). For \(w\in W\),
\(e(x+N_gx,w)=e(x,w)\) gives \(e(N_gx,w)=1\), so \(N_g(A[2])\) lies in the line \(W^\perp\).
Hence \(N_g\) is one of the two maps \(A[2]/W\to W^\perp\), and the image of \(G_{\Q}\) in
\(\GL(A[2])\) has order at most \(2\).

Let \(E\subset A_K\) be an elliptic curve defined over \(K\).
\begin{enumerate}
\item \emph{\(E[2]\neq E^\sigma[2]\) as subgroups of \(A\).} Every lift \(\tilde\sigma\) maps
\(E[2]\) onto \(E^\sigma[2]\neq E[2]\), so it acts nontrivially on \(A[2]\), as the unique
nontrivial element \(u\) of the image. If some \(h\in G_K\) acted as \(u\), the lift
\(\tilde\sigma h\) would act trivially. So \(G_K\) acts trivially on \(A[2]\), and
\(E[2]\subseteq E(K)\). Take \(E_1=E\).

\item \emph{\(E[2]=E^\sigma[2]=:\Lambda\).} The subgroup \(\Lambda\) is stable under \(G_K\) and
under \(\tilde\sigma\), hence under \(G_{\Q}\). Write \(\iota:E\to A_K\) for the inclusion,
\(e_2^E:E[2]\times E^\vee[2]\to\mu_2\) for the canonical pairing and \(\phi_E:E\to E^\vee\) for the
principal polarization, so that \(e_E(x,y)=e_2^E(x,\phi_E y)\) is the Weil pairing of \(E\). By
functoriality, \(e(\iota x,z)=e_2^E(x,\iota^\vee\lambda z)\) for \(x\in E[2]\) and \(z\in A[2]\).
The polarization \(\iota^\vee\lambda\iota\) of \(E\) is \(n\phi_E\) for some \(n\ge1\), so
\(e(\iota x,\iota y)=e_E(x,y)^n\) for \(x,y\in E[2]\).

If \(n\) is even, \(\Lambda\) is isotropic, hence \(\Lambda=\Lambda^\perp\), and \(x\mapsto
e(x,\cdot)\) identifies \(A[2]/\Lambda\) with the dual \(\Lambda^\vee\) as \(G_{\Q}\)-modules. If
\(\Lambda\not\subseteq W\), then \(W\cap\Lambda\) is a line and \(W+\Lambda=A[2]\), so \(G_{\Q}\)
acts trivially on \(A[2]/\Lambda\cong\Lambda^\vee\), hence on \(\Lambda\), and
\(\Lambda\subseteq W\) after all. So \(E[2]=\Lambda\subseteq A(\Q)\), and \(E[2]\subseteq E(K)\).
Take \(E_1=E\).

If \(n\) is odd, \(e\) is nondegenerate on \(\Lambda\), so \(A[2]=\Lambda\oplus\Lambda^\perp\) as
\(G_{\Q}\)-modules and \(W=\Lambda^{G_{\Q}}\oplus(\Lambda^\perp)^{G_{\Q}}\). As \(\dim W=3\),
either \(\Lambda\subseteq W\), and we take \(E_1=E\), or \(\Lambda^\perp\subseteq W\). In the latter
case let \(E^\perp\) be the complementary elliptic curve, the connected component of the kernel of
\(\iota^\vee\lambda:A_K\to E^\vee\); it is defined over \(K\). For \(x\in E[2]\) and \(y\in
E^\perp[2]\) we have \(e(\iota x,y)=e_2^E(x,\iota^\vee\lambda y)=1\), so
\(E^\perp[2]\subseteq\Lambda^\perp\). Both groups have order \(4\), so
\(E^\perp[2]=\Lambda^\perp\subseteq A(\Q)\), and \(E^\perp[2]\subseteq E^\perp(K)\). Finally
\((E^\perp)^\sigma\neq E^\perp\): otherwise \(E^\perp\) would be stable under \(G_{\Q}\), hence
defined over \(\Q\), contradicting the simplicity of \(A\) over \(\Q\). Take \(E_1=E^\perp\).
\end{enumerate}
In all cases \(E_1\) has the required properties. \qed

\begin{statement}{Lemma 5 (Weil restriction).}\label{lemma-res}
Let \(E\subset A_K\) be an elliptic curve defined over \(K\) with \(E^\sigma\neq E\).
\begin{enumerate}
\item[(a)] \(A\) is isogenous over \(\Q\) to \(\Res_{K/\Q}E'\) for every elliptic curve \(E'\) over
\(K\) that is \(K\)-isogenous to \(E\).
\item[(b)] Let \(p\) be an odd prime, inert in \(K\), at which \(A\) has good reduction, and put
\(\mathfrak p=pO_K\). Then every such \(E'\) has good reduction at \(\mathfrak p\), and
\(\#\widetilde A(\F_p)=\#\widetilde{E'}(O_K/\mathfrak p)\), where \(O_K/\mathfrak p\cong\F_{p^2}\).
\item[(c)] If \(A(\Q)\) has a point of order \(5\) and \(E[2]\subseteq E(K)\), then some \(E'\)
that is \(K\)-isogenous to \(E\) has \(E'(K)\supseteq\Z/2\times\Z/10\).
\end{enumerate}
\end{statement}

\emph{Proof.} (a) Let \(\iota:E\to A_K\) be the inclusion and \(\iota^\sigma:E^\sigma\to A_K\) its
conjugate. The map \((x,y)\mapsto\iota(x)+\iota^\sigma(y)\) from \((\Res_{K/\Q}E)(\bar\Q)\) to
\(A(\bar\Q)\) is a homomorphism of abelian varieties over \(\bar\Q\), and it commutes with the
action of \(G_{\Q}\) described in (I7), because \(g\circ\iota=\iota^g\circ g\) for \(g\in G_{\Q}\).
By Galois descent it is a homomorphism \(\Res_{K/\Q}E\to A\) over \(\Q\). It is surjective because
\(E+E^\sigma=A_{\bar\Q}\), hence an isogeny. A \(K\)-isogeny \(E'\to E\) induces a \(\Q\)-isogeny
\(\Res_{K/\Q}E'\to\Res_{K/\Q}E\).

(b) Fix a prime \(\ell\neq p\). By (a) and (I7), \(V_\ell(A)\cong\Ind_{G_K}^{G_{\Q}}V_\ell(E')\)
as \(G_{\Q}\)-modules. By (I1)(a) this is unramified at \(p\). As \(p\) is unramified in \(K\),
the inertia groups at primes above \(p\) lie in \(G_K\), and \(V_\ell(E')\) is a
\(G_K\)-submodule of \(V_\ell(A)\); so it is unramified at \(\mathfrak p\), and \(E'\) has good
reduction at \(\mathfrak p\). Let \(g\in G_{\Q}\) be an arithmetic Frobenius element at a prime
above \(p\), well defined up to conjugation and inertia, which acts trivially. It does not lie in
\(G_K\), as \(p\) is inert, and \(g^2\in G_K\) is an arithmetic Frobenius element at
\(\mathfrak p\), because \(O_K/\mathfrak p\) has degree \(2\) over \(\F_p\). In a decomposition
\(\Ind V=V\oplus gV\), \(g\) has the block form \(\left(\begin{smallmatrix}0&B\\C
&0\end{smallmatrix}\right)\), so \(\det(1-g)=\det(1-BC)\), where \(BC\) is the action of
\(g^2\) on \(V_\ell(E')\). By (I1),
\(\#\widetilde A(\F_p)=\det(1-g\mid V_\ell(A))=\det(1-g^2\mid V_\ell(E'))
=\#\widetilde{E'}(O_K/\mathfrak p)\).

(c) Let \(P_5\in A(\Q)\) have order \(5\) and let \(\pi:A_K\to Y=A_K/E^\sigma\) be the quotient
map; its restriction to \(E\) is a \(K\)-isogeny \(E\to Y\). By (I3) over \(K\), either
\(P_5\in E^\sigma\), and then \(P_5=\tilde\sigma(P_5)\in E(K)\), or \(\pi(P_5)\in Y(K)\) has order
\(5\). So there is a \(K\)-isogeny \(f:E\to E_5\) with \(E_5\in\{E,Y\}\) and a point of order \(5\)
in \(E_5(K)\). Let \(H\) be the \(5\)-primary part of \(\ker f\), \(a:E\to E'=E/H\) and \(b:E'\to
E_5\) with \(f=b\circ a\). As \(\deg a\) is a power of \(5\), \(a\) maps \(E[2]\) isomorphically
and \(G_K\)-equivariantly onto \(E'[2]\), so \(E'[2]\subseteq E'(K)\). As \(\deg b\) is prime to
\(5\), the dual isogeny \(\hat b\) is injective on \(E_5[5]\) and maps the point of order \(5\)
to a point of order \(5\) in \(E'(K)\), as in the proof of \cite[Lemma~1]{P63}. \qed

\begin{statement}{Lemma 6 (a Chebotarev criterion).}\label{lemma-cheb}
Let \(K=\Q(\sqrt d)\) be a quadratic field and \(E:y^2=(x-e_1)(x-e_2)(x-e_3)\) an elliptic curve
over \(K\) with \(e_1,e_2,e_3\in K\), and put \(N_{ij}=N_{K/\Q}(e_i-e_j)\in\Q^\times\). Suppose
that \(8\) divides \(\#\widetilde E(O_K/pO_K)\) for all but finitely many primes \(p\) inert in
\(K\). Then, for a suitable numbering \(\{i,j,k\}=\{1,2,3\}\),
\[N_{ij}\in\Q^{\times2}\qquad\text{and}\qquad N_{ik}\in\Q^{\times2}\ \text{ or }\ N_{jk}\in\Q^{\times2}
\ \text{ or }\ d\,N_{ik}N_{jk}\in\Q^{\times2}.\]
\end{statement}

\emph{Proof.} For all but finitely many inert \(p\), we have \(p>3\), the \(e_i\) are integral
at \(\mathfrak p=pO_K\) with unit differences, and the reduction
\(\widetilde E\) over \(k=O_K/\mathfrak p\cong\F_{p^2}\) has full rational \(2\)-torsion. Then
\(8\mid\#\widetilde E(k)\) if and only if \(\widetilde E(k)\) contains a point of order \(4\),
that is, some \((\bar e_m,0)\) lies in \(2\widetilde E(k)\); by (I8) this means that
\(\bar e_m-\bar e_{m'}\) is a square in \(k\) for both \(m'\neq m\). An element of
\(\F_{p^2}^\times\) is a square if and only if its norm to \(\F_p\) is, and since \(\sigma\)
induces the Frobenius of \(k/\F_p\), the norm of \(\bar e_m-\bar e_{m'}\) is \(N_{mm'}\bmod p\).
So for all but finitely many inert \(p\):
\[(\ast)\qquad\text{for some }m,\ \Bigl(\frac{N_{mm'}}{p}\Bigr)=1\text{ for both }m'\neq m .\]
Let \(V\) be the subgroup of \(\Q^\times/\Q^{\times2}\) generated by the classes of \(d\),
\(N_{12}\), \(N_{13}\), \(N_{23}\), written additively as an \(\F_2\)-vector space. By Chebotarev's
theorem for \(\Q(\sqrt d,\sqrt{N_{12}},\sqrt{N_{13}},\sqrt{N_{23}})\), every homomorphism
\(\phi:V\to\F_2\) is \(x\mapsto\bigl(1-(x/p)\bigr)/2\) for infinitely many primes \(p\), and these
\(p\) are inert in \(K\) when \(\phi(d)=1\). So \((\ast)\) says: no \(\phi\) with \(\phi(d)=1\)
takes the value \(1\) on two of \(N_{12}\), \(N_{13}\), \(N_{23}\). A system \(\phi(v_1)=\dots=
\phi(v_s)=1\) is solvable unless some sum of an odd number of the \(v_t\) vanishes in \(V\). For
each pair \(\{x,y\}\subset\{N_{12},N_{13},N_{23}\}\) the system \(\phi(d)=\phi(x)=\phi(y)=1\) is
therefore unsolvable, so \(x=0\), \(y=0\) or \(x+y=d\) in \(V\) (\(d\neq0\) as \(K\neq\Q\)). If
none of the three classes were \(0\), then \(N_{12}+N_{13}=N_{12}+N_{23}=N_{13}+N_{23}=d\), so
all three would be equal and \(d=0\). Hence some \(N_{ij}\) is a square, and the pair
\(\{N_{ik},N_{jk}\}\) gives the second condition. \qed

\medskip
The condition is also sufficient: if it holds, then \((\ast)\) holds at every inert \(p\) as in
the proof (if \(dN_{ik}N_{jk}\) is a square, exactly one of \(N_{ik}\), \(N_{jk}\) is a square
modulo \(p\)), so \(8\) divides \(\#\widetilde E(O_K/pO_K)\). The curve
\(y^2=x(x-1)(x-\sqrt2)\) over \(\Q(\sqrt2)\), with \(N_{12}=1\), \(N_{13}=-2\), \(N_{23}=-1\) and
\(dN_{13}N_{23}=4\), shows that the third alternative is needed.

\medskip
Let \(\mathcal X\) be the elliptic curve \(y^2=x^3-2x^2+1\) over \(\Q\), with origin the point at
infinity \(O\) (Cremona \(20a2\)). Put
\[
e_1=-\frac{t^3(t-1)}{t^2-3t+1},\quad
B_t=\frac{-8t^4+12t^3-4t+1}{(t^2-3t+1)^2},\quad
\rho=\frac{(2t-1)^2s}{(t^2-3t+1)^2},\quad
e_{2,3}=\frac{-B_t\pm\rho}{8}.
\]

\begin{statement}{Lemma 7 (\(\Z/2\times\Z/10\) over quadratic fields).}\label{lemma-x110}
\begin{enumerate}
\item[(a)] \(\mathcal X\) has rank \(0\) and \(\mathcal X(\Q)=\{O,(0,\pm1),(1,0),(2,\pm1)\}\).
\item[(b)] Let \(K\) be a quadratic field and \(E'\) an elliptic curve over \(K\) with
\(E'(K)\supseteq\Z/2\times\Z/10\). Then \(E'\cong E_t\) over \(K\) for some \(t\in K\) such that
\(s^2=(2t-1)(4t^2-2t-1)\) for some \(s\in K\), and \(E'\) is isomorphic over \(K\) to
\(y^2=(x-e_1)(x-e_2)(x-e_3)\). The point \(Q=(2t,s)\in\mathcal X(K)\) does not lie in
\(\mathcal X(\Q)\).
\item[(c)] Put \(R=Q+Q^\sigma\in\mathcal X(\Q)\). If \(R=O\), then \(t\in\Q\). If \(R\neq O\), then
the line through \(Q\) and \(Q^\sigma\) passes through \(-R\) and has a slope \(m\in\Q\); the
\(x\)-coordinates of \(Q\) and \(Q^\sigma\) are the roots of the monic quadratic
\[q_{R,m}(X)=\frac{X^3-2X^2+1-\bigl(y(-R)+m(X-x(-R))\bigr)^2}{X-x(-R)}\in\Q[X],\]
which is irreducible over \(\Q\), and \(K=\Q\bigl(\sqrt{\disc q_{R,m}}\bigr)\).
\end{enumerate}
\end{statement}

\emph{Proof.} (a) is part B of the certificate (Magma proves the rank is \(0\); the torsion is
\(\Z/6\)). (b) By (I9), \((E',P)\cong(E_t,(0,0))\) for a point \(P\) of order \(10\); as \(E_t\) is
an elliptic curve, \(t\neq\tfrac12\) and \(t^2-3t+1\neq0\). The certificate (part C) checks that
\(e_1=x(5\cdot(0,0))\) is a root of the \(2\)-division polynomial \(4x^3+b_2x^2+2b_4x+b_6\) of
\(E_t\), that the complementary quadratic factor is \(4x^2+B_tx+C_t\) with
\[B_t^2-16C_t=\frac{(2t-1)^5(4t^2-2t-1)}{(t^2-3t+1)^4}
=(2t-1)(4t^2-2t-1)\cdot\Bigl(\frac{(2t-1)^2}{(t^2-3t+1)^2}\Bigr)^2,\]
and that over \(\Q(t)(s)\), \(s^2=(2t-1)(4t^2-2t-1)\), the \(2\)-division polynomial is
\(4\prod(x-e_i)\). Since \(E'[2]\subseteq E'(K)\), this discriminant is a square in \(K\); the
factor \((2t-1)^2/(t^2-3t+1)^2\) is nonzero, so \((2t-1)(4t^2-2t-1)=s^2\) with \(s\in K\), and
completing the square identifies \(E_t\) with \(y^2=\prod(x-e_i)\). The six
points of \(\mathcal X(\Q)\) have \(t=x/2\in\{\infty,0,\tfrac12,1\}\), where \(E_t\) is not
defined or is singular (certificate, part B); also no elliptic curve over \(\Q\) has a subgroup
\(\Z/2\times\Z/10\) \cite{Ma}. (c) The point \(R\) is fixed by \(\sigma\), so \(R\in\mathcal X(\Q)\).
If \(R=O\), then \(Q^\sigma=-Q\) has the same \(x\)-coordinate, so \(x(Q)\in\Q\). If \(R\neq O\),
then \(Q\neq Q^\sigma\) and the line through them is \(\sigma\)-stable, hence defined over \(\Q\);
its third intersection with \(\mathcal X\) is \(-(Q+Q^\sigma)=-R\). It is not vertical, since the
vertical line through \(-R\) meets \(\mathcal X\) only in \(\pm R\) and \(O\). The remaining
statements follow. \qed

\begin{statement}{Proposition 8.}\label{prop-families}
Let \(R\in\mathcal X(\Q)\setminus\{O\}\) and \(m\in\Q\) be such that \(q_{R,m}\) is irreducible
over \(\Q\). Let \(x_Q\) be a root of \(q_{R,m}\), \(K=\Q(x_Q)\), \(t=x_Q/2\) and
\(s=y(-R)+m(x_Q-x(-R))\), so that \(Q=(x_Q,s)\in\mathcal X(K)\), and suppose that \(E_t\) is an
elliptic curve. Then the conclusion of \hyperref[lemma-cheb]{Lemma~6} fails for
\(y^2=(x-e_1)(x-e_2)(x-e_3)\) over \(K\), with \(e_1,e_2,e_3\) as above.
\end{statement}

\emph{Proof.} The other root of \(q_{R,m}\) gives \(Q^\sigma\) and the conjugates of the \(e_i\),
with the same norms, so the choice of \(x_Q\) does not matter. Over the function field \(\Q(m)\),
let \(K_m=\Q(m)[X]/(q_{R,m})\),
\(t=X/2\), \(s=y(-R)+m(X-x(-R))\); then \(s^2=(2t-1)(4t^2-2t-1)\), and the norms
\(N_{ij}\in\Q(m)\) of \(e_i-e_j\) and \(d=\disc q_{R,m}\) are computed in part D of the
certificate. Up to squares in \(\Q(m)^\times\), in each of the five families two of the three
norms have a common class \(g_R\), and the third is \(-1\) or \(-d\):
\[
\begin{array}{c|c|c|c}
R & d=\disc q_{R,m} & g_R & \text{third norm}\\\hline
(0,1) & m^4+4m^2-8m+4 & (m-1)(m^2+m-1) & -1\\
(0,-1) & m^4+4m^2+8m+4 & -(m+1)(m^2-m-1) & -1\\
(2,-1) & m^4-8m^2+8m & -(m-1)(m^2-3m+1) & -d\\
(2,1) & m^4-8m^2-8m & (m+1)(m^2+3m+1) & -d\\
(1,0) & m^4-2m^2+5 & (m^2-1)(m^2-5) & -1
\end{array}
\]
Let \(m_0\in\Q\) be as in the statement, and suppose first that \(m_0\) is not a root of any
numerator or denominator of \(d\), of the \(N_{ij}\), or of \(N_{K_m/\Q(m)}(t^2-3t+1)\). Let
\(\mathcal L=\Q[m]_{(m-m_0)}\) be the local ring of \(\Q[m]\) at \(m_0\). As \(q_{R,m}\) is monic in \(X\) with
coefficients in \(\Q[m]\), the ring \(\mathcal O=\mathcal L[X]/(q_{R,m})\) is free over
\(\mathcal L\) with basis \(1,X\), and \(\mathcal O/(m-m_0)\mathcal O\cong\Q[X]/(q_{R,m_0})\cong K\),
with \(X\mapsto x_Q\). Norms from \(\mathcal O\) to \(\mathcal L\) are determinants of
multiplication maps in this basis, so they commute with reduction modulo \(m-m_0\). The element
\(t^2-3t+1\) of \(\mathcal O\) has norm \(N_{K_m/\Q(m)}(t^2-3t+1)\in\Q[m]\), which does not vanish
at \(m_0\), so it is a unit of \(\mathcal O\). The \(e_i\) have denominators dividing
\(8(t^2-3t+1)^2\), so they lie in \(\mathcal O\), and their reductions are the \(e_i\) of the
statement. Hence \(N_{ij}(m_0)=N_{K/\Q}(e_i-e_j)\), and \(d(m_0)=\disc q_{R,m_0}\). Each identity
of square classes in the table reads \(N=\kappa h^2\) with \(\kappa\in\{g_R,-1,-d\}\) and
\(h\in\Q(m)^\times\). The zeros and poles of \(h\) lie among those of \(N\), \(g_R\) and \(d\), and
the roots of \(g_R\) are roots of numerators or denominators of the \(N_{ij}\). So
\(h(m_0)\in\Q^\times\), \(g_R(m_0)\neq0\), and the classes above hold at \(m_0\). The conclusion of \hyperref[lemma-cheb]{Lemma~6} would need some
norm to be a square. The class \(-1\) is not a square. If the third norm is \(-d\) and is a
square, the second condition requires \(g_R(m_0)\in\Q^{\times2}\) or \(d\cdot g_R(m_0)^2\in
\Q^{\times2}\); the latter would make both \(d\) and \(-d\) squares. So in every case
\(g_R(m_0)\in\Q^{\times2}\), and \((m_0,y)\) is a rational point on \(C_R:y^2=g_R(m)\). The five
curves \(C_R\) have genus \(1\) and Jacobians \(40a3\) (for \(R=(0,\pm1),(1,0)\)) and \(20a2\) (for
\(R=(2,\pm1)\)), of rank \(0\) (part E of the certificate). Their rational points are the points
at infinity and those with \(m\in\{0,1\}\), \(\{-1,0\}\), \(\{0,1,2\}\), \(\{-2,-1,0\}\),
\(\{\pm1\}\) for \(R=(0,1),(0,-1),(2,-1),(2,1),(1,0)\) respectively; for each of these \(m\),
\(q_{R,m}\) splits over \(\Q\). Finally, part F of the certificate treats the excluded values
\(m_0\) directly: each gives a reducible \(q_{R,m_0}\), a value of \(t\) with \(t^2-3t+1=0\), or a
singular \(E_t\). \qed

\medskip
\textbf{The quadratic case.} Let \(A\) be principally polarized, simple over \(\Q\), of type (N)
or of type (S) with \(\Gamma\cong C_2\), with a point of order \(5\) in \(A(\Q)\), and suppose
\(A(\Q)[2]\) has rank at least \(3\). By \hyperref[lemma-full]{Lemma~4} there is
\(E_1\subset A_K\) over \(K\) with \(E_1^\sigma\neq E_1\) and \(E_1[2]\subseteq E_1(K)\). By
\hyperref[lemma-res]{Lemma~5}(c) some \(E'\), \(K\)-isogenous to \(E_1\), has
\(E'(K)\supseteq\Z/2\times\Z/10\), and by Lemma~5(a) \(A\) is isogenous over \(\Q\) to
\(\Res_{K/\Q}E'\). The subgroup of \(A(\Q)\) generated by the point of order \(5\) and
\(A(\Q)[2]\) has order divisible by \(40\); by (I1)(b) and Lemma~5(b),
\(40\mid\#\widetilde A(\F_p)=\#\widetilde{E'}(O_K/pO_K)\) for every odd prime \(p\neq5\) inert in
\(K\) at which \(A\) has good reduction. Let \(E'\cong E_t\) and \(Q\in\mathcal X(K)\) be as in
\hyperref[lemma-x110]{Lemma~7}.
\begin{itemize}
\item If \(Q+Q^\sigma=O\), then \(t\in\Q\), and \(E'\) is the base change of the elliptic curve
\(E_t\) over \(\Q\). By (I7), \(\Res_{K/\Q}E'\) is isogenous over \(\Q\) to
\(E_t\times E_t^{(K)}\), so \(A\) is not simple over \(\Q\). This is a contradiction.
\item Otherwise \(R=Q+Q^\sigma\neq O\), and by \hyperref[lemma-x110]{Lemma~7}(c) the line through
\(Q\) and \(Q^\sigma\) has a slope \(m\in\Q\), \(q_{R,m}\) is irreducible with root \(x(Q)\),
\(K=\Q(x(Q))\), and \(y(Q)=y(-R)+m(x(Q)-x(-R))\). So \(R\), \(m\) and \(x_Q=x(Q)\) satisfy the
hypotheses of \hyperref[prop-families]{Proposition~8}. The curve \(E'\) is isomorphic over \(K\)
to \(y^2=\prod(x-e_i)\), so the two have isomorphic reductions at all but finitely many primes
of \(K\), and \hyperref[lemma-cheb]{Lemma~6} applies to \(y^2=\prod(x-e_i)\). Its conclusion
contradicts \hyperref[prop-families]{Proposition~8}.
\end{itemize}
Hence \(A(\Q)[2]\) has rank at most \(2\).

\section{Proof of the Theorem and consequences for the census}\label{census}

\emph{Proof of the Theorem.} Let \(A/\Q\) be principally polarized, simple over \(\Q\) and
geometrically split, with a rational point of order \(5\). By (I2), \(A\) is of type (N), (CM) or
(S). In type (CM), \(A(\Q)[2]\) has rank at most \(2\) by Section~\ref{cm}. In type (S),
\(\Gamma\cong C_2\) or \(C_4\) by \hyperref[prop-gamma]{Proposition~2}; if \(\Gamma\cong C_4\),
then \(\#A(\Q)[2^\infty]\le4\) by \hyperref[prop-A]{Proposition~A}. In type (N) and in type (S)
with \(\Gamma\cong C_2\), the rank is at most \(2\) by Section~\ref{quadratic}. \qed

\medskip
\emph{Proof of the Corollary.} The Jacobian of a curve of genus \(2\) over \(\Q\) is a principally
polarized abelian surface, and each of \([2,2,10]\), \([2,2,20]\) and \([2,2,2,10]\) contains a
point of order \(5\) and has \(2\)-rank at least \(3\). For Jacobians split over \(\Q\), the
\href{https://genus-2-torsion.github.io/proofs/qsplit-two-rank-obstruction.html}{\(2\)-rank
obstruction} shows that rational \(2\)-rank \(4\) excludes rational points of order \(5\). \qed

\medskip
On the census, the three groups become impossible in the class ``geometrically split, simple
over \(\Q\)''. For \([2,2,2,10]\), already impossible in the class ``split over \(\Q\)'', the
combined geometrically split cell becomes impossible (\(\emptyset\)). The group \([2,2,10]\) has
examples split over \(\Q\): the census records the curve 2.2.10.b and a certified infinite family
of Jacobians split over \(\Q\) with exact torsion group \([2,2,10]\) (see its
\href{https://genus-2-torsion.github.io/group.html?g=2.2.10}{census page} and the
\href{https://genus-2-torsion.github.io/data/knowledge/sources/family_2210_split_certificate.log}{family
certificate}). The class ``split over \(\Q\)'' of \([2,2,20]\) stays open. So the combined cells
of \([2,2,10]\) and \([2,2,20]\) keep their grades.

\section{Remarks}\label{remarks}

\begin{enumerate}
\item \emph{The point of order \(5\) is needed} in types (N), (CM) and (S) with
\(\Gamma\in\{C_2,V_4,D_8\}\): the LMFDB curves 2560.a.819200.1, 4096.b.65536.1, 25600.f.512000.1,
8192.b.131072.1 and 12800.c.128000.1 are simple over \(\Q\), geometrically split, with rational
torsion subgroup \((\Z/2)^3\) and Sato--Tate groups \(N(\mathrm{SU}(2)\times\mathrm{SU}(2))\),
\(J(C_2)\), \(J(E_1)\), \(J(E_2)\) and \(J(E_4)\) (component groups \(C_2\), \(D_2\), \(C_2\),
\(D_2\), \(D_4\)). For \(\Gamma\cong C_4\) (Sato--Tate group \(E_4\)) no point of order \(5\) is
needed. Consistently, the LMFDB curves with Sato--Tate group \(E_4\) that are simple over \(\Q\)
have torsion \([2]\), \([2,2]\) or \([2,10]\); the last is \(X_1(16)\) (256.a.512.1), whose
endomorphism field is \(\Q(\zeta_{16})^+\), and for the primes \(p<200\) with \(p\equiv3,11
\pmod{16}\) one finds \(\#J_1(16)(\F_p)\equiv4\pmod8\) (for example \(\#J_1(16)(\F_3)=20=2^2+4^2\),
\(u=-1\) in the proof of Proposition~A).

\item \emph{A check of Lemmas~4--6 on known examples (not used in the proof).} For the eight
LMFDB curves with rational torsion subgroup \((\Z/2)^3\) that are simple over \(\Q\) and split
over a quadratic field \(K\) (2169.a.175689.1, 2457.b.199017.1, 2560.a.5120.1, 2560.a.819200.1,
4096.b.65536.1, 4352.b.278528.1, 8192.b.131072.1, 25600.f.512000.1), the supplementary script
\texttt{qsimple\_2210\_validation.m} checks that \(\#J(\F_p)=\#E(\F_{p^2})\) at all inert primes
\(p<400\) of good reduction for the LMFDB elliptic factor \(E\) over \(K\) (and the product formula
at split primes), finds in the \(2\)-isogeny class of \(E\) over \(K\) a curve with full
\(K\)-rational \(2\)-torsion, as \hyperref[lemma-full]{Lemma~4} predicts, and checks that each such
curve satisfies the conclusion of \hyperref[lemma-cheb]{Lemma~6}, as it must since
\(8\mid\#J(\F_p)\). The curve 4096.b.65536.1 has CM and 8192.b.131072.1 has \(\Gamma\cong V_4\),
so they are not of the kind treated in Section~\ref{quadratic}; but in the proofs of Lemmas~4
and~5 the type of \(A\) enters only through the simplicity of \(A\) over \(\Q\) and an elliptic
curve \(E\subset A_K\) defined over \(K\), and Lemma~6 concerns elliptic curves over \(K\). The script also checks the example
after Lemma~6: \(8\) divides \(\#\widetilde E(\F_{p^2})\) at every inert prime \(p<400\).

\item \emph{The polarization.} Only \hyperref[lemma-full]{Lemma~4} uses that the polarization is
principal. Section~\ref{cm}, \hyperref[prop-gamma]{Proposition~2} and
\hyperref[prop-A]{Proposition~A} hold for every abelian surface over \(\Q\) of the stated types.

\item \emph{Why the argument for order \(7\) does not apply.} The proof for \([2,2,14]\) used
that no elliptic curve over a quadratic field has a subgroup \(\Z/28\) or \(\Z/2\times\Z/14\)
\cite[Lemma~1]{P2214}. Over quadratic
fields \(\Z/2\times\Z/10\) does occur, for infinitely many fields; the quadratic points of
\(X_1(2,10)\) are therefore controlled here by the Frobenius condition at inert primes.
\end{enumerate}

\section{Computations}\label{computations}

The certificate
\href{https://genus-2-torsion.github.io/data/knowledge/sources/qsimple_2210_checks.m}{\texttt{qsimple\_2210\_checks.m}}
(Magma V2.30-1, under a second) consists of assertions and ends with \texttt{ALL\_CHECKS\_PASSED};
its log is
\href{https://genus-2-torsion.github.io/data/knowledge/sources/qsimple_2210_checks.log}{\texttt{qsimple\_2210\_checks.log}}.
SHA-256 of the script:
\texttt{b5ee08ba8354b767\allowbreak 816e762666e18aa2\allowbreak 177c000fd03fd2b8\allowbreak 97f8d6d86f01bc05}; of the log:
\texttt{29c3bc4a93d80343\allowbreak 76ee575ce865fecb\allowbreak ff3ebcbabce9e834\allowbreak 70a1e0dd9ff070ee}. Part A lists the values
\(h(4)\) of \hyperref[lemma-f3]{Lemma~1}. Part B checks that \(\mathcal X\) is \(20a2\) with rank
\(0\) (proved by Magma) and torsion \(\Z/6\), lists \(\mathcal X(\Q)\), and checks that \(E_t\) is
singular at \(t=0,\tfrac12,1\). Part C checks the Kubert family, the formula \(t=c^2/(b-c)\), the
discriminant in \hyperref[lemma-x110]{Lemma~7}(b) and the formulas for \(e_1,e_2,e_3\) over
\(\Q(t)(s)\). Part D computes, for each \(R\neq O\), the quadratic \(q_{R,m}\), the field
\(K_m\), the identity \(s^2=(2t-1)(4t^2-2t-1)\), the identity
\(4x^3+b_2x^2+2b_4x+b_6=4\prod(x-e_i)\) for the specialized Kubert curve, the norms \(N_{ij}\) and
their square classes, and asserts the table of \hyperref[prop-families]{Proposition~8},
including the discriminants \(d\). Part E proves that each \(C_R\) has rank \(0\), checks the
Cremona label of its Jacobian, checks that the points found are all of \(C_R(\Q)\) (their number
equals the order of the torsion of the Jacobian) and that their finite slopes are those listed in
Proposition~8, and checks that each of these slopes gives a reducible \(q_{R,m}\). Part F computes the excluded slopes of Proposition~8 (rational roots of all
numerators and denominators involved) and checks from scratch that each gives a reducible
\(q_{R,m}\), a value of \(t\) with \(t^2-3t+1=0\), or a singular \(E_t\). Part G checks, at the
sample slope \(m=7/3\) in each family, that \(Q+Q^\sigma=R\), that \(E_t(K)_{\mathrm{tors}}\cong
\Z/2\times\Z/10\), and that the conclusion of Lemma~6 fails there.

The supplementary script
\href{https://genus-2-torsion.github.io/data/knowledge/sources/qsimple_2210_validation.m}{\texttt{qsimple\_2210\_validation.m}}
of Remark~2 (a few seconds) also consists of assertions and ends with \texttt{VAL\_DONE}; its log
is
\href{https://genus-2-torsion.github.io/data/knowledge/sources/qsimple_2210_validation.log}{\texttt{qsimple\_2210\_validation.log}}.

\begin{thebibliography}{P2214}

\bibitem[Kn]{Kn}
A. W. Knapp, \emph{Elliptic Curves}, Mathematical Notes~40, Princeton University Press, 1992:
Theorem~4.2 (pp.~85--87).
\href{https://doi.org/10.1515/9780691186900}{doi:10.1515/9780691186900}.

\bibitem[Ku]{Ku}
D. S. Kubert, \emph{Universal bounds on the torsion of elliptic curves}, Proc. London Math. Soc.
(3) \textbf{33} (1976), 193--237: Table~3.
\href{https://doi.org/10.1112/plms/s3-33.2.193}{doi:10.1112/plms/s3-33.2.193}.

\bibitem[Ma]{Ma}
B. Mazur, \emph{Modular curves and the Eisenstein ideal}, Publ. Math. IH\'ES \textbf{47} (1977),
33--186. \href{https://doi.org/10.1007/BF02684339}{doi:10.1007/BF02684339}.

\bibitem[Mi]{Mi}
J. S. Milne, \emph{On the arithmetic of abelian varieties}, Invent. Math. \textbf{17} (1972),
177--190: \S1. \href{https://doi.org/10.1007/BF01425446}{doi:10.1007/BF01425446}.

\bibitem[P63]{P63}
Census proof of 21 October 2026, \emph{An obstruction to rational \(63\)-torsion on geometrically
split abelian surfaces},
\url{https://genus-2-torsion.github.io/proofs/qsimple-63-obstruction.html}: Section~1,
Lemmas~1--3 and~5--7, Sections~5--7 and Section~9.1.

\bibitem[P2214]{P2214}
Census proof of 10 November 2026, \emph{Rational \(7\)-torsion and the \(2\)-rank of
geometrically split abelian surfaces},
\url{https://genus-2-torsion.github.io/proofs/qsimple-2214-obstruction.html}: Section~1,
Section~2 (types (N) and (S) with \(\Gamma=C_2\)), Section~3 and Lemmas~1 and~3.

\bibitem[ST]{ST}
J.-P. Serre and J. Tate, \emph{Good reduction of abelian varieties}, Ann. of Math. (2)
\textbf{88} (1968), 492--517: Theorem~1 and Corollary~2 (p.~493), Lemma~2 (p.~495), \S2
(pp.~496 and~498). \href{https://doi.org/10.2307/1970722}{doi:10.2307/1970722}.

\end{thebibliography}

\end{document}
